% Figure: a two-leg sling lifting a load. (a) the rigging scene with the
% included angle and the leg angle from vertical; (b) the leg-tension
% multiplier as a function of the angle each leg makes with the vertical.
% Numbers match the multi-leg-sling case study in the text.
\begin{figure}[htbp]
\centering
\begin{tikzpicture}[
  leg/.style={draw=engblue, line width=1.4pt},
  ten/.style={draw=engred, -{Stealth}, very thick},
  wt/.style={draw=engblack, -{Stealth}, very thick},
  scale=1.0,
]
% ===== Left: the rigging scene =====
\begin{scope}[xshift=0cm]
  % hook / master link at top
  \coordinate (H) at (2.0,4.0);
  \draw[draw=engblack, very thick] (H) -- ++(0,0.5);
  \draw[draw=engblack, fill=white, thick] (2.0,4.65) circle (0.18);
  % attachment points on the load
  \coordinate (L) at (0.6,1.4);
  \coordinate (R) at (3.4,1.4);
  % sling legs
  \draw[leg] (H) -- (L);
  \draw[leg] (H) -- (R);
  % the load (a crate)
  \draw[draw=engblack, fill=engblue!10, thick] (0.4,0.4) rectangle (3.6,1.4);
  \node[font=\scriptsize] at (2.0,0.9) {load $W$};
  \fill[engblack] (L) circle (0.06);
  \fill[engblack] (R) circle (0.06);
  % vertical reference through the hook
  \draw[draw=enggray, thin, dashed] (H) -- ++(0,-2.6);
  % leg angle from vertical at the hook
  \draw[draw=enggray, thin] (2.0,3.3) arc (270:236:0.7);
  \node[font=\scriptsize] at (1.55,3.25) {$\theta$};
  % included angle marker
  \draw[draw=enggray, thin] (1.3,3.1) arc (236:304:0.7);
  \node[font=\scriptsize] at (2.0,2.95) {$2\theta$};
  \node[font=\small, anchor=north] at (2.0,0.1) {(a) two-leg sling};
\end{scope}
% ===== Right: tension multiplier vs leg angle =====
\begin{scope}[xshift=6.6cm, yshift=0.2cm]
  \begin{axis}[
    width=6.6cm, height=5.4cm,
    xlabel={leg angle from vertical $\theta$ (deg)},
    ylabel={leg tension $/$ (half load)},
    xmin=0, xmax=75, ymin=0.9, ymax=4.2,
    xtick={0,15,30,45,60,75},
    ytick={1,2,3,4},
    grid=both, grid style={enggray!25},
    label style={font=\scriptsize},
    tick label style={font=\scriptsize},
    every axis plot/.append style={line width=1.2pt},
  ]
    % multiplier = 1/cos(theta)
    \addplot[draw=engred, domain=0:74, samples=80] {1/cos(x)};
    \addplot[draw=enggray, dashed, domain=0:75] {2};
    \node[font=\scriptsize, anchor=west, engred] at (axis cs:6,1.7) {$1/\cos\theta$};
    \node[font=\scriptsize, anchor=south east, enggray] at (axis cs:60,2.0) {tension $= W$};
  \end{axis}
  \node[font=\small, anchor=north] at (3.0,-0.4) {(b) tension penalty};
\end{scope}
\end{tikzpicture}
\caption[A two-leg sling and the tension penalty of a wide angle]{(a) A
symmetric two-leg sling lifts a load $W$; each leg makes an angle $\theta$
with the vertical, so the legs subtend an included angle $2\theta$ at the
hook. Vertical balance of the master link gives $2T\cos\theta = W$, so each
leg carries $T = W/(2\cos\theta)$. (b) The leg tension, in units of half
the load, rises as $1/\cos\theta$: at $\theta = 0$ each leg carries exactly
half the load; at $\theta = 60^{\circ}$ each leg carries the full load $W$;
beyond that the penalty climbs without bound. The horizontal pull that the
wide angle introduces is the silent cost a rigger pays for spreading the
legs to clear the load, and it is why sling charts derate sharply below a
$45^{\circ}$ included angle.}
\label{fig:vol03ch01:sling-plan}
\end{figure}
